\chapter{Общий Хиггс и запрет прямой межрёберной кривизны второй степени}

\section{Постановка после интуитивного аудита}

Полный гессиан исправленного минимума оставил двадцать семь нулевых мод,
поскольку три физические кромки имеют независимые семейные кадры. Перед
вычислением полного универсального исчисления необходимо проверить более
сильное и дешёвое условие: существует ли вообще ненулевое прямое
смешанное слово второй степени на уже принятом физическом носителе с одним
общим хиггсовским дублетом.

Ранняя семейная ветвь проекта уже проверяла отдельные операторы
расщепления. Факторные трансляции различают три направления, но не выводят
свои относительные веса и имеют общий собственный базис; естественные
симметрийные усреднения возвращают разложение $1+2$. Поэтому настоящий
гейт не повторяет этот путь и не вставляет новый семейный оператор.

\section{Три физические кромки при общем дублете}

Пусть $H\in\mathbb C^2$, $\|H\|=1$, и
\begin{equation}
 \widetilde H=i\sigma_2\overline H.
 \label{eq:v7-common-higgs-conjugate}
\end{equation}
Тогда
\begin{equation}
 H^\dagger\widetilde H=0,
 \qquad
 \|\widetilde H\|=1.
 \label{eq:v7-higgs-orthogonal-pair}
\end{equation}

Обозначим через
\begin{equation}
 T_u,T_d,T_e:H_L^8\longrightarrow H_R^7
 \label{eq:v7-physical-edge-maps}
\end{equation}
три разрешённые карты: $T_u$ сворачивает каждый цветной слабый дублет с
$\widetilde H$, $T_d$ --- с $H$, а $T_e$ сворачивает лептонный дублет с
$H$. Их правые концы $u_R,d_R,e_R$ попарно ортогональны.

Из ортогональности слабых направлений и правых концов следует
\begin{equation}
 T_aT_b^\dagger=0,
 \qquad
 T_a^\dagger T_b=0,
 \qquad a\ne b.
 \label{eq:v7-physical-edge-degree-two-orthogonality}
\end{equation}
Для пары $u,d$ первое равенство использует
$\widetilde H^\dagger H=0$, а второе --- ортогональность $u_R$ и $d_R$.
Для пар с $e$ дополнительно различаются кварковый и лептонный левые блоки.

\section{Подъём на исправленный аффинный носитель}

Пусть
\begin{equation}
 X_a\in E_{\rm aff}=\operatorname{Hom}(\mathbb C^4,\operatorname{im}P_3),
 \qquad
 D_a=X_a\otimes T_a.
 \label{eq:v7-lifted-edge-arrows}
\end{equation}
Тензорное произведение сохраняет~\eqref{eq:v7-physical-edge-degree-two-orthogonality}:
\begin{equation}
 D_aD_b^\dagger=0,
 \qquad
 D_a^\dagger D_b=0,
 \qquad a\ne b.
 \label{eq:v7-lifted-degree-two-cross-zero}
\end{equation}
Следовательно, для $D=D_u+D_d+D_e$
\begin{align}
 D^\dagger D&=\sum_aD_a^\dagger D_a,\label{eq:v7-source-square-sum}\\
 DD^\dagger&=\sum_aD_aD_a^\dagger.\label{eq:v7-target-square-sum}
\end{align}
То же верно для Hodge-коммутатора нечётных стрелок:
\begin{equation}
 [d,d^\dagger]=\sum_a[d_a,d_a^\dagger].
 \label{eq:v7-hodge-commutator-edge-sum}
\end{equation}

Тем самым прямой квадрат текущего оператора содержит только собственные
матрицы Грама трёх кромок. Он не содержит $X_aX_b^\dagger$ или
$X_a^\dagger X_b$ с $a\ne b$.

\section{Физическая бимодульная проекция и junk}

Предыдущая классификация одноформ дала точную матрицу размерностей
межрёберных интертвинеров
\begin{equation}
 \left(\dim\operatorname{Hom}_{\mathcal A_{\rm SM}-\mathcal A_{\rm SM}}
 (\lambda_a,\lambda_b)\right)_{a,b}
 =I_3.
 \label{eq:v7-edge-intertwiner-identity}
\end{equation}
Поэтому нуль-форменные вставки могут создавать сырые операторные пути, но
их внедиагональные части не являются эндоморфизмами физического
бимодуля $\lambda_u\oplus\lambda_d\oplus\lambda_e$. После физической
проекции размерность смешанного угла равна нулю.

Факторизация по junk-идеалу является факторизацией уже представленного
пространства форм. Она может удалить существующий класс, но не создать
смешанный класс, который равен нулю до факторизации и отсутствует в
физическом эндоморфизмном углу.

\begin{theorem}[Прямой запрет второй степени]
На исправленном носителе
$E_{\rm aff}\otimes\Lambda_{\rm ch}$ при одном общем нормированном
хиггсовском дублете все прямые смешанные произведения второй степени между
кромками $u,d,e$ равны нулю. Физическая бимодульная проекция также не
содержит внедиагональных межрёберных эндоморфизмов. Поэтому обычная
физическая кривизна второй степени остаётся суммой трёх независимых
кромочных кривизн и не поднимает относительные $U(3)$-моды.
\end{theorem}

\begin{remark}[Точная область утверждения]
Теорема не объявляет классифицированным всё сырое универсальное исчисление
с произвольными нуль-форменными вставками. Она закрывает более узкий
физический маршрут: прямую common-Higgs кривизну второй степени после
проекции на эндоморфизмы уже принятого трёхрёберного бимодуля.
\end{remark}

\section{Граница более высокой степени}

Спектральное кручение на $H_{15}$ уже содержит смешанную чувствительность
\begin{equation}
 I_{ud}=3|y_u|^2|y_d|^2,
 \label{eq:v7-torsion-higher-degree-witness}
\end{equation}
но прежний гейт доказал, что оно измеряет заданные юкавские отношения, а
не выбирает их без нового вариационного принципа. Тем не менее эта формула
показывает правильную степень следующего поиска: первая возможная
межрёберная информация возникает не в прямом квадрате, а в выведенном
квартичном или более высоком циклическом слове.

Такой член допускается только если он:
\begin{enumerate}
 \item следует из того же Hodge/суперсвязностного родителя;
 \item имеет фиксированную нормировку общего следа;
 \item корректно типизирован до взятия следа;
 \item остаётся ненулевым после физической и junk-проекций;
 \item не использует наблюдаемые массы или смешивания как вход.
\end{enumerate}

\begin{nogocriterion}[Запрет ручного портала после нулевой двухформы]
Нельзя заменять исчезнувший класс второй степени членом
$\|X_u-X_d\|^2$, назначенной матрицей смешивания или произвольной нормой
компонент кручения. После настоящего no-go допустим только независимо
выведенный инвариант более высокой степени.
\end{nogocriterion}

\section{Вердикт}

\begin{protocol}
\begin{enumerate}
 \item общий нормированный хиггсовский дублет: \textbf{зафиксирован};
 \item точное равенство $H^\dagger\widetilde H=0$: \textbf{получено};
 \item прямые смешанные слова $T_aT_b^\dagger$: \textbf{ноль};
 \item прямые смешанные слова $T_a^\dagger T_b$: \textbf{ноль};
 \item сохранение нуля после аффинного подъёма: \textbf{получено};
 \item разложение $D^\dagger D$ и $DD^\dagger$ по кромкам:
 \textbf{точное};
 \item смешанные физические бимодульные эндоморфизмы:
 \textbf{отсутствуют};
 \item junk может породить отсутствующий класс: \textbf{нет};
 \item полный сырой универсальный калькулюс: \textbf{не заявлен};
 \item подъём двадцати семи относительных нулей второй степенью:
 \textbf{закрыт отрицательно};
 \item следующий гейт: \textbf{допуск коэффициент-свободного
 квартичного межрёберного инварианта}.
\end{enumerate}
\end{protocol}

Машинный сертификат сохранён в
\path{s2t_v7_common_higgs_degree_two_cross_edge_gate.py} и
\path{s2t_v7_common_higgs_degree_two_cross_edge_gate_results.json}.